\section{Filtered \texorpdfstring{\mbox{\normalsize$\jepPubScript{D}$}}{D}-modules associated to \texorpdfstring{\mbox{\normalsize$\jepPubBold{Q}$}}{Q}-divisors}\label{coverings}
Let $X$ be a smooth complex algebraic variety, with $\dim(X)=n$. The ideals we associate to effective $\jepPubBold{Q}$-divisors on $X$ arise from certain Hodge $\jepPubScript{D}$-modules. 
The $\jepPubScript{D}$-modules themselves have been extensively studied: these are the 
$\jepPubScript{D}$-modules attached to rational powers of functions on $X$. We proceed to recall their definition.

Consider a nonzero $h\in\jepPubScript{O}_X(X)$ and $\beta\in\jepPubBold{Q}$. We denote by $Z$ the reduced divisor on $X$ with the same support as
$H={\rm div}(h)$ and let $j\colon U=X\smallsetminus {\rm Supp}(Z)\hookrightarrow X$ be the inclusion map. 
 We consider the left $\jepPubScript{D}_X$-module 
$\jepPubScript{M}(h^{\beta})$, which is a rank 1 free $\jepPubScript{O}_X(*Z)$-module 
with generator the symbol $h^{\beta}$, on which a derivation $D$ of $\jepPubScript{O}_X$ acts by
$$D(wh^{\beta}) :=\Big(D(w) +  w\frac{\beta\cdot D(h)}{h}\Big)h^{\beta}.$$
We will denote the corresponding right $\jepPubScript{D}_X$-module by $\jepPubScript{M}_r(h^{\beta})$.
This can be described as $h^{\beta}\omega_X(*Z)$, an $\jepPubScript{O}_X$-module isomorphic to $\omega_X(*Z)$,
and such that if $x_1,\ldots,x_n$ are local coordinates, then 
$$(h^{\beta}w dx_1\cdots dx_n)\partial_i=- h^{\beta}\Big(\frac{\partial w}{\partial x_i}+ w\frac{\beta}{h}\cdot \frac{\partial h}{\partial x_i}\Big)dx_1\cdots dx_n$$
for every $i$ with $1\leqslant i\leqslant n$. 



\begin{remark}\label{case_ell_1}
When $\beta\in\jepPubBold{Z}$, we have a canonical isomorphism of left $\jepPubScript{D}_X$-modules
\begin{equation}\label{eq_case_ell_1}
\jepPubScript{M}(h^{\beta})\simeq \jepPubScript{O}_X(*Z),\,\,\,\,\,\,wh^{\beta}\longmapsto w \cdot h^{\beta},
\end{equation}
where on the localization $\jepPubScript{O}_X(*Z)$ we consider the natural $\jepPubScript{D}_X$-module structure induced from $\jepPubScript{O}_X$. Note that $\jepPubScript{O}_X(*Z)$
is also the $\jepPubScript{D}$-module push-forward $j_+ \jepPubScript{O}_U$.
\end{remark}

\begin{remark}\label{eq_renormalization_0}
For every positive integer $m$, we have a canonical isomorphism of  left $\jepPubScript{D}_X$-modules
$$\jepPubScript{M}(h^{\beta})\simeq\jepPubScript{M}\big((h^m)^{\beta /m}\big),\,\,\,\,\,\,wh^{\beta}\longmapsto w(h^m)^{\beta/m}.$$
\end{remark}


\begin{remark}\label{twist_individual}
We can define, more generally, left $\jepPubScript{D}$-modules $\jepPubScript{M}(h_1^{\beta_1}\cdots h_r^{\beta_r})$, for 
nonzero regular functions $h_1,\ldots,h_r\in\jepPubScript{O}_X(X)$ and rational numbers $\beta_1,\ldots,\beta_r$. 
If $\ell_i$ are positive integers such that
$\beta_i/\ell_i=\beta$ for all $i$ and if $h=\prod_ih_i^{\ell_i}$, then we have an isomorphism of left $\jepPubScript{D}_X$-modules
$$\jepPubScript{M}(h_1^{\beta_1}\cdots h_r^{\beta_r})\simeq \jepPubScript{M}(h^{\beta}).$$
\end{remark}

\begin{remark}\label{twist_positive}
If $r$ is an integer, then we have an isomorphism of left $\jepPubScript{D}_X$-modules
$$\jepPubScript{M}(h^{\beta})\longrightarrow \jepPubScript{M}(h^{r + \beta}), \,\,\,\,\,\,wh^{\beta}\longmapsto (wh^{-r}) h^{r+\beta}.$$
\end{remark}

\medskip

Let now $D$ be an effective $\jepPubBold{Q}$-divisor on $X$. We denote by $Z$ the reduced divisor with the same support as $D$. 
As above, we put $U=X\smallsetminus Z$ and let $j\colon U\hookrightarrow X$ be the inclusion map. 
We first assume that we can write $D=\alpha\cdot {\rm div}(h)$ for some nonzero $h\in\jepPubScript{O}_X(X)$ and $\alpha\in\jepPubBold{Q}_{>0}$ (this is of course 
always the case locally). To this data we can associate the $\jepPubScript{D}_X$-module $\jepPubScript{M}(h^{-\alpha})$; later it will be more convenient to 
consider equivalently (according to Remark \ref{twist_positive}) the $\jepPubScript{D}_X$-module $\jepPubScript{M}(h^{1-\alpha})$. 
This depends on the choice of $h$; however, if we replace $h$ by $h^m$ and $\alpha$ by $\alpha/m$, for some positive integer $m$,
the $\jepPubScript{D}$-module does not change (see Remark~\ref{eq_renormalization_0}). In particular, we may always assume that $\alpha=1/\ell$,
for a positive integer $\ell$.




\begin{remark}\label{rmk_add_eff_div}
Suppose that $D'$ is a $\jepPubBold{Q}$-divisor with the same support as $D$ and such that $D-D'={\rm div}(u)$, for some $u\in\jepPubScript{O}_X(X)$.
Suppose that we can write
$D'=(1/\ell)\cdot {\rm div}(h')$ for some $h'\in\jepPubScript{O}_X(X)$ and some positive integer $\ell$.
In this case we can also write
$D=(1/\ell)\cdot {\rm div}(h)$, where $h=u^{\ell}{h'}$, and  we
have an isomorphism of  $\jepPubScript{D}_X$-modules
\begin{equation}\label{isom_add_1}
\jepPubScript{M}(h^{-1/\ell})\longrightarrow \jepPubScript{M}({h'}^{-1/\ell}), \,\,\,\,\,\,gh^{-1/\ell} \longmapsto gu^{-1}{h'}^{-1/\ell}.
\end{equation}
\end{remark}

\medskip

Our first goal is to show that $\jepPubScript{M}(h^{-\alpha})$ is canonically a filtered $\jepPubScript{D}_X$-module. 
Let $\ell$ be a positive integer such that $\ell\alpha\in\jepPubBold{Z}$.
Consider the finite \'{e}tale map $p\colon V\to U$, where $V={\bf Spec}\,\jepPubScript{O}_U[y]/(y^{\ell}-h^{-\ell\alpha})$.
Note that this fits in a Cartesian diagram
\begin{equation}\label{cart_diag1}
{\def\labelstyle{\textstyle}\xymatrix{V\ar[r]\ar[d]_{p}&W\ar[d]^{q}\\U\ar[r]_{j}&X,}}
\end{equation}
in which
$$W={\bf Spec}\,\jepPubScript{O}_X[z]/(z^\ell-h^{\ell\alpha}),$$
such that the map $V\to W$ pulls $z$ back to $y^{-1}=y^{\ell-1}h^{\ell\alpha}$. 


\begin{lemma}\label{lem_decomp1}
We have an isomorphism of left $\jepPubScript{D}_X$-modules
\begin{equation}\label{eq1_lem_decomp1}
j_+p_+\jepPubScript{O}_V\simeq\mathop{\textstyle\bigoplus}\limits_{i=0}^{\ell-1}\jepPubScript{M}(h^{-i\alpha}),
\end{equation}
with the convention that the first summand is $\jepPubScript{O}_X(*Z)$.
\end{lemma}

\begin{proof}
Since $p$ is finite \'{e}tale, it follows that we have a canonical isomorphism $\tau\colon p^*\jepPubScript{D}_U\simeq\jepPubScript{D}_V$, 
and for every $\jepPubScript{D}_V$-module $\jepPubScript{M}$ we have $p_+\jepPubScript{M}\simeq p_*\jepPubScript{M}$,
with the action of $\jepPubScript{D}_U$ induced via the isomorphism $\tau$.

By mapping $gy^i$ to $gh^{-i\alpha}$, where $g$ is a section of $\jepPubScript{O}_X$ and $0\leqslant i\leqslant\ell-1$, we obtain an isomorphism 
of $\jepPubScript{O}_X$-modules as in
(\ref{eq1_lem_decomp1}). In order to see that this is an isomorphism of $\jepPubScript{D}_X$-modules, consider a local derivation
$D$ of $\jepPubScript{O}_X$ and note that since $y^{\ell}=h^{-\ell\alpha}$, by identifying $D$ with its pull-back to $V$ we have
$$D (y^i)=iy^{i-1}D(y)=-i\alpha y^i\frac{D(h)}{h},$$
which via our map corresponds to $D (h^{-i\alpha})$.
This implies the assertion.
\end{proof}

It follows from the lemma that the right-hand side of (\ref{eq1_lem_decomp1}) is the $\jepPubScript{D}$-module corresponding to the 
mixed Hodge module push-forward $(j\circ p)_+\jepPubBold{Q}_V^H[n]$. In particular, it 
 carries a canonical structure of filtered $\jepPubScript{D}$-module. 


\begin{remark}\label{rem_replace_by_multiple}
Let us see what happens if we replace $\ell$ by a multiple $m\ell$. Let $p_{\ell}\colon V_{\ell}\to U$ and $p_{m\ell}\colon V_{m\ell}\to U$ be the corresponding \'{e}tale covers.
Note that $$V_{m\ell}={\bf Spec}\,\jepPubScript{O}_U[y]/(y^{\ell m}-h^{-\ell m\alpha})$$ decomposes as a disjoint union of $m$ copies of $V_{\ell}$, and thus we have an isomorphism of filtered $\jepPubScript{D}_X$-modules (and a corresponding isomorphism of mixed Hodge modules)
\begin{equation}\label{eq_rem_replace_by_multiple}
j_+(p_{m\ell})_+\jepPubScript{O}_{V_{m\ell}}\simeq 
\big(j_+(p_{\ell})_+\jepPubScript{O}_{V_{\ell}}\big)^{\oplus m}.
\end{equation}
If $\eta$ is a primitive root of $1$ of order $\ell m$, and if on each side of 
(\ref{eq_rem_replace_by_multiple}) we consider the decompositions 
(\ref{eq1_lem_decomp1}), then the isomorphism maps 
$$h^{-i\alpha} \longmapsto \big(\eta^{is}h^{-c\ell\alpha}\cdot h^{-d\alpha}\big)_{0\leqslant s\leqslant m-1},$$
where we write $i=\ell c+d$, with $0\leqslant c\leqslant m-1$ and $0\leqslant d\leqslant\ell-1$.
\end{remark}

We can interpret the isomorphism in (\ref{eq1_lem_decomp1}) in terms of a suitable $\mu_\ell$-action, where $\mu_{\ell}$
is the group of $\ell$-th roots of $1$ in $\jepPubBold{C}^*$. Note that we have a natural action of $\mu_{\ell}$ on $W$ such that
via the corresponding action on $\jepPubScript{O}_W$, an element $\lambda\in\mu_{\ell}$ maps $z^i$ to $\lambda^iz^i$. 
If we let $\mu_{\ell}$ act trivially on $X$, then $q$ is an equivariant morphism (in fact, $q$ is the quotient morphism with respect
to the $\mu_{\ell}$-action). It is clear that $q^{-1}(Z)$ is fixed by the $\mu_{\ell}$-action and we have an induced $\mu_{\ell}$-action
on $W\smallsetminus q^{-1}(Z)=V$. This in turn induces a $\mu_{\ell}$-action on $j_+p_+\jepPubScript{O}_V$ and the isomorphism in
(\ref{eq1_lem_decomp1}) corresponds to the isotypic decomposition of $j_+p_+\jepPubScript{O}_V$, such that every 
$\lambda\in\mu_{\ell}$ acts on $\jepPubScript{M}(h^{-i\alpha})$ by multiplication with $\lambda^{-i}$. 


\begin{lemma}\label{decomp_filtration}
The filtration on $j_+p_+\jepPubScript{O}_V$ is preserved by the $\mu_{\ell}$-action. Therefore we have an induced filtration on each 
$\jepPubScript{M}(h^{-i\alpha})$ such that (\ref{eq1_lem_decomp1}) is an isomorphism of filtered $\jepPubScript{D}$-modules.
\end{lemma}


\begin{proof}
One way to see this is by using a suitable equivariant resolution of $W$. Let $W'$ be the disjoint union of the irreducible components of $W$ and
$q'\colon W'\to W$ the canonical morphism. It is clear that the $\mu_{\ell}$-action on $W$ induces an action on $W'$ such that $q'$ is equivariant.
Since $V$ is contained in the smooth locus of $W$, it has an open immersion into $W'$. 
We use equivariant resolution of singularities to construct a $\mu_{\ell}$-equivariant morphism 
$\varphi\colon Y\to W'$ that is an isomorphism over $V$ and such that $(q\circ q'\circ\varphi)^*(Z)$ is a divisor with simple normal crossings.  
Let $g=q\circ q'\circ \varphi$. If $E$ is the reduced, effective divisor supported on $g^{-1}(Z)$, then 
we have an isomorphism of filtered $\jepPubScript{D}$-modules (induced by a corresponding isomorphism of mixed Hodge modules)
\begin{equation}\label{eq_ism_filtered}
j_+p_+\jepPubScript{O}_V\simeq g_+\widetilde{j}_+\jepPubScript{O}_V\simeq g_+\jepPubScript{O}_Y(*E),
\end{equation}
where $\widetilde{j}\colon Y\smallsetminus {\rm Supp}(E)\hookrightarrow Y$ is the inclusion map.

We can deduce the assertion in the lemma from an explicit
computation of the filtration on $j_+p_+\jepPubScript{O}_V$ via the isomorphism (\ref{eq_ism_filtered}), as follows. First, since we deal with
$\jepPubScript{D}$-module push-forward, it is more convenient to work with right $\jepPubScript{D}$-modules. We will thus compute $g_+\omega_Y(*E)$, where
$\omega_Y(*E)$ is the filtered right $\jepPubScript{D}$-module corresponding to $\jepPubScript{O}_Y(*E)$. 

Since $E$ is a simple normal crossing divisor, $\omega_Y(*E)$ has a resolution by a complex $C^\jepPubBullet$ 
of filtered right $\jepPubScript{D}_Y$-modules
$$0\longrightarrow C^{-n}\longrightarrow\cdots\longrightarrow C^0\longrightarrow 0,$$
where $C^{i}=\Omega_Y^{i+n}(\log E)\otimes_{\jepPubScript{O}_Y}\jepPubScript{D}_Y$, with the filtration given by
$$F_{k-n}C^i=\Omega_Y^{i+n}(\log E)\otimes_{\jepPubScript{O}_Y}F_{k+i}\jepPubScript{D}_Y.$$ 
For a description of the maps in this complex, see the beginning of Section~\ref{complex_for_SNC} below;
a proof
of the fact that it resolves $\omega_Y(*E)$ is given in \cite[Prop.~3.1]{MP1}. We can thus compute  $F_kg_+\omega_Y(*E)$ as the image of the injective map
$$R^0g_*\big(F_k(C^{\jepPubBullet}\otimes_{\jepPubScript{D}_Y}\jepPubScript{D}_{Y\to X})\big)\longrightarrow R^0g_*(C^{\jepPubBullet}\otimes_{\jepPubScript{D}_Y}\jepPubScript{D}_{Y\to X})=g_+\omega_Y(*E).$$
Since $g$ is equivariant and the action of $\mu_{\ell}$ on $Y$ induces an action on $E$ (in fact, it fixes $E$), the above description implies that each
$F_kg_+\omega_Y(*E)$ is preserved by the $\mu_{\ell}$-action.
\end{proof}

\begin{remark}\label{rmk_filtration_first_piece}
We note that the filtration on $j_+p_+\jepPubScript{O}_V$ induces the canonical filtration on the first summand $\jepPubScript{O}_X(*Z)$.
Indeed, on $U$ we have a morphism of mixed Hodge modules $\jepPubBold{Q}_U^H[n]\to p_+\jepPubBold{Q}_V^H[n]$. Applying $j_+$ and only considering the underlying filtered $\jepPubScript{D}$-modules, we obtain a morphism $j_+\jepPubScript{O}_U\to j_+p_+\jepPubScript{O}_V$, which is an isomorphism onto the first summand.
\end{remark}

\begin{definition}
Given $\alpha>0$, choose $\ell\geqslant 2$ such that $\ell\alpha\in\jepPubBold{Z}$. In this case $\jepPubScript{M}(h^{-\alpha})$ appears as the second summand in the decomposition
(\ref{eq1_lem_decomp1}). We define the filtration 
$$F_k \jepPubScript{M}(h^{-\alpha}) \,\,\,\,\,\,\text{for} \,\,\,\,\,\, k \geqslant 0$$ 
to be the filtration induced from the canonical filtration on $j_+p_+\jepPubScript{O}_V$. 
It is straightforward to see, using the discussion in Remark~\ref{rem_replace_by_multiple} that this filtration does not change if we replace $\ell$
by a multiple; therefore it is independent of $\ell$. Moreover, we note that if $\alpha$ is an integer, using the same Remark~\ref{rem_replace_by_multiple}, 
the isomorphism $\jepPubScript{M}(h^{-\alpha})\simeq\jepPubScript{O}_X(*Z)$ is an isomorphism of filtered $\jepPubScript{D}$-modules.
\end{definition}

In this definition, a priori different covers have to be considered for each of the summands $\jepPubScript{M}(h^{-i\alpha})$. However, 
we have:

\begin{lemma}\label{lem_decomp1_v2}
With the filtration defined above, the isomorphism (\ref{eq1_lem_decomp1}) is an isomorphism of filtered $\jepPubScript{D}$-modules. 
\end{lemma}
\begin{proof}
By Lemma~\ref{decomp_filtration}, we only need to show that for every $i$ with $0\leqslant i\leqslant \ell-1$, the filtration induced on $\jepPubScript{M}(h^{-i\alpha})$ by that on $j_+p_+\jepPubScript{O}_V$ coincides with the 
one given in the above definition. For $i=0$ this follows from Remark~\ref{rmk_filtration_first_piece}. If $i>0$, consider the cover used to define the filtration on 
$\jepPubScript{M}(h^{-i\alpha})$, namely 
$$p'\colon V'={\bf Spec}\jepPubScript{O}_U[y]/(y^{\ell}-h^{i\alpha\ell})\longrightarrow U.$$ 
Note that we have a finite morphism
$\psi\colon V\to V'$ of varieties over $U$, that pulls-back $y$ to $y^i$. We have a canonical morphism of mixed Hodge modules
$\jepPubBold{Q}_{V'}^H[n]\to \psi_+\jepPubBold{Q}_V^H[n]$. Applying $j_+p'_+$ and passing to the underlying filtered $\jepPubScript{D}$-modules, we obtain a morphism of filtered 
$\jepPubScript{D}$-modules $j_+p'_+\jepPubScript{O}_{V'}\to j_+p_+\jepPubScript{O}_V$ that is the identity on the summand $\jepPubScript{M}(h^{-i\alpha})$. This proves our claim.
\end{proof}

\begin{remark}\label{same_filtration_rescaling}
It is clear from definition that for every $\alpha>0$ and every positive integer $m$, the isomorphism
$$\jepPubScript{M}(h^{-\alpha})\longrightarrow \jepPubScript{M}\big((h^m)^{-\alpha/m}\big), \quad gh^m\longmapsto g (h^m)^{-\alpha/m}$$
is an isomorphism of filtered $\jepPubScript{D}$-modules. 
\end{remark}


\begin{remark}\label{rem_isom_comp_with_filtrations}
In the setting of Remark~\ref{rmk_add_eff_div},
the isomorphism (\ref{isom_add_1}) is an isomorphism of filtered $\jepPubScript{D}_X$-modules.
This is clear if $\ell=1$, hence we assume $\ell\geqslant 2$.
Let
$p\colon V\to U$ and $p'\colon V'\to U$ be the canonical projections, where
$$V={\bf Spec}\,\jepPubScript{O}_U[y]/(y^{\ell}-h)\quad\text{and}
\quad V'={\bf Spec}\,\jepPubScript{O}_U[y]/(y^{\ell}-h').$$
We have an isomorphism $\varphi\colon V'\to V$ of schemes over $U$, where $\varphi^*(y)=uy$. 
This induces an isomorphism of filtered $\jepPubScript{D}_X$-modules
$$j_+p_+\jepPubScript{O}_V\simeq j_+p'_+\jepPubScript{O}_{V'},$$
which via the identifications given by Lemma~\ref{lem_decomp1} is the direct sum
$$\mathop{\textstyle\bigoplus}\limits_{i=0}^{\ell-1}\jepPubScript{M}(h^{-i/\ell})\simeq 
\mathop{\textstyle\bigoplus}\limits_{i=0}^{\ell-1}\jepPubScript{M}({h'}^{-i/\ell})$$
of the isomorphisms (\ref{isom_add_1}). For $i=1$, we obtain our assertion.
\end{remark}

A special case of the above remark implies that for every $\alpha>0$ the isomorphism
$$\jepPubScript{M}(h^{-\alpha})\longrightarrow\jepPubScript{M}(h^{-\alpha-1}),\quad gh^{-\alpha}\longmapsto (gh)h^{-\alpha-1}$$
is an isomorphism of filtered $\jepPubScript{D}$-modules. We use this to put a structure of filtered
$\jepPubScript{D}$-module on $\jepPubScript{M}(h^{\beta})$ for every $\beta\in\jepPubBold{Q}$, such that for every $r\in\jepPubBold{Z}$,
we have an isomorphism of filtered $\jepPubScript{D}$-modules
$$\jepPubScript{M}(h^{\beta})\longrightarrow\jepPubScript{M}(h^{\beta-r}),\quad gh^{\beta}\longmapsto (gh^r)h^{\beta-r}.$$
For example, we have have an isomorphism of filtered $\jepPubScript{D}$-modules $\jepPubScript{M}(h^0)\simeq \jepPubScript{O}_X(*Z)$.


\begin{remark}\label{rem_indep_equations1}
Suppose that $h, \overline{h} \in\jepPubScript{O}_X(X)$ are nonzero, and $\alpha,\overline{\alpha}\in\jepPubBold{Q}_{>0}$ are 
such that we have the equality of $\jepPubBold{Q}$-divisors
$$\alpha \cdot {\rm div}(h)=\overline{\alpha}\cdot {\rm div}(\overline{h}).$$
Let $\ell$ be  a positive integer such that $\ell\alpha,\ell\overline{\alpha}\in\jepPubBold{Z}$.
In this case there is $g\in\jepPubScript{O}_X^*(X)$ such that $h^{\ell\alpha}=g{\overline{h}}^{\ell\overline{\alpha}}$.
 Suppose now that there exists $G\in\jepPubScript{O}_X(X)$ such that $G^{\ell}=g$. (For example, this holds after pulling-back to the
\'{e}tale cover $\mathbf{Spec}\,\jepPubScript{O}_X[z]/(z^{\ell}-g)$.)
In this case we have an isomorphism of filtered $\jepPubScript{D}_X$-modules
$$\Phi\colon \jepPubScript{M}(h^{-\alpha})\longrightarrow \jepPubScript{M}(\overline{h}^{-\overline{\alpha}})$$
given by
$$\Phi(wh^{-\alpha})=wG^{-1}\overline{h}^{-\overline{\alpha}}.$$
Indeed, this follows from the definition of the filtrations and the isomorphism of schemes over $U$
$$\varphi\colon {\bf Spec} \jepPubScript{O}_U[y]/(y^{\ell}-\overline{h}^{-\ell\overline{\alpha}})\longrightarrow {\bf Spec} \jepPubScript{O}_U[y]/(y^{\ell}-h^{-\ell\alpha})$$
that pulls-back $y$ to $G^{-1}y$. 
\end{remark}

\begin{remark}\label{rem_behavior_etale}
It is clear that the filtration on $\jepPubScript{M}(h^{-\alpha})$ is compatible with restriction to open subsets.
More generally, it is compatible with smooth pullback, as follows.
Suppose that $h\in\jepPubScript{O}_X(X)$ is nonzero and $\alpha \in \jepPubBold{Q}$.
If $\varphi\colon Y\to X$ is a smooth morphism and $g=h\circ \varphi$, then there is an isomorphism of $\jepPubScript{D}_Y$-modules
$$\jepPubScript{M}(g^{-\alpha})\simeq\varphi^*\jepPubScript{M}(h^{-\alpha}),$$
such that for every $k$ we have
$$F_k\jepPubScript{M}(g^{-\alpha})\simeq\varphi^*F_k\jepPubScript{M}(h^{-\alpha}).$$
Indeed, choose $\ell\geqslant 2$ such that $\ell\alpha\in\jepPubBold{Z}$ and consider the 
Cartesian diagram
$$
{\def\labelstyle{\textstyle}\xymatrix{V_Y\ar[r]^{p_Y}\ar[d]_{\psi}&U_Y\ar[r]^{j_Y}\ar[d]&Y\ar[d]^{\textstyle\varphi}\\V\ar[r]_{p}&U\ar[r]_{j}&X,}}
$$
where $j$ and $p$ are as in Lemma~\ref{lem_decomp1} and $j_Y$ and $p_Y$ are the corresponding morphisms for $Y$ and $g$. 
Note that we have a base-change theorem that gives
\begin{equation}\label{eq_rem_behavior_etale}
\varphi^!j_+p_+\jepPubBold{Q}^H_V[n]\simeq (j_Y)_+(p_Y)_+\psi^!\jepPubBold{Q}^H_V[n]
\end{equation}
(see \cite[(4.4.3)]{Saito-MHM}).
Moreover, since $\varphi$ is smooth, if $d=\dim(Y)-\dim(X)$, then for every filtered $\jepPubScript{D}$-module 
$(\jepPubScript{M},F)$ underlying a mixed Hodge module $M$, the filtered $\jepPubScript{D}$-module
underlying $\varphi^!M$ is $(\varphi^*\jepPubScript{M},F)[d]$, where $F_k(\varphi^*\jepPubScript{M})=\varphi^*(F_{k}\jepPubScript{M})$
(see \cite[3.5]{Saito-MHP}).
This also applies to $\psi$; in particular, we have $\psi^!\jepPubBold{Q}_V^H[n]\simeq\jepPubBold{Q}_{V_Y}^H[n+2d]$.
By decomposing both sides of (\ref{eq_rem_behavior_etale}) with respect to the $\mu_{\ell}$-action, we
obtain our assertion.
\end{remark}










